\section{Arsitektur Sistem yang Diusulkan}

Arsitektur sistem yang diusulkan dalam penelitian ini menggunakan pendekatan ABAC berbasis \textit{edge} pada lingkungan IoT. Arsitektur ini dirancang untuk mendekatkan proses evaluasi keputusan akses ke \textit{gateway edge}, sehingga \textit{endpoint} IoT tidak perlu menjalankan proses otorisasi yang kompleks dan tidak seluruh permintaan akses harus selalu bergantung pada \textit{cloud}. Secara umum, sistem terdiri atas empat komponen utama, yaitu \textit{endpoint IoT, gateway edge, proxy/server}, dan \textit{cloud/PIP}.

\textit{Endpoint IoT} berperan sebagai penghasil permintaan akses. Setiap \textit{request} membawa informasi yang dibutuhkan dalam evaluasi ABAC, seperti atribut subjek, objek, aksi, dan kondisi lingkungan. \textit{Request} tersebut dikirimkan ke \textit{gateway edge} untuk diproses. Pada penelitian ini, \textit{endpoint} tidak menjalankan evaluasi kebijakan secara langsung, melainkan hanya mengirimkan permintaan akses kepada \textit{gateway}. Hal ini dilakukan agar beban evaluasi kebijakan yang lebih kompleks dapat dipindahkan ke \textit{node edge} yang memiliki sumber daya lebih besar dibanding \textit{endpoint} IoT.

\textit{Gateway edge} merupakan komponen utama dalam arsitektur yang diusulkan. Pada \textit{gateway} ditempatkan PEP, PDP, dan \textit{cache} atribut lokal. PEP menerima \textit{request} dari \textit{endpoint} dan meneruskannya ke PDP untuk dievaluasi. PDP bertugas mengevaluasi kebijakan ABAC berdasarkan atribut yang dibutuhkan. Setelah PDP menghasilkan keputusan akses, PEP menegakkan keputusan tersebut dalam bentuk \textit{allow} atau \textit{deny}, kemudian mengirimkan \textit{decision response} kembali ke \textit{endpoint}.

\textit{Cache} atribut lokal pada \textit{gateway} digunakan untuk menyimpan atribut tertentu yang dibutuhkan dalam evaluasi kebijakan. \textit{Cache} ini bersifat lokal per \textit{gateway}, sehingga atribut yang tersimpan pada satu \textit{gateway} tidak dibagikan secara langsung ke \textit{gateway} lain. Pendekatan ini digunakan untuk merepresentasikan kondisi \textit{edge} yang terdistribusi, di mana setiap \textit{gateway} memiliki \textit{state cache} masing-masing. Penggunaan \textit{cache} bertujuan untuk mengurangi jumlah pengambilan atribut dari PIP, menurunkan latensi evaluasi akses, dan mengurangi \textit{overhead} komunikasi ke \textit{cloud}.

\textit{Cloud/PIP} berperan sebagai sumber atribut aktual atau \textit{ground truth}. Apabila atribut tidak tersedia di \textit{cache}, TTL atribut telah kedaluwarsa, atau atribut tersebut tidak dikonfigurasi sebagai atribut yang boleh di-\textit{cache}, maka \textit{gateway} akan mengirimkan \textit{attribute query} ke \textit{cloud/PIP} melalui \textit{proxy/server}. \textit{Cloud/PIP} kemudian mengembalikan nilai atribut aktual pada \textit{gateway}. Selain menyediakan atribut aktual, \textit{cloud/PIP} juga merepresentasikan sumber pembaruan atribut, termasuk perubahan status perangkat, perubahan \textit{role}, perubahan kondisi \textit{trust}, dan \textit{revocation update}.

\textit{Proxy/server} ditempatkan di antara \textit{gateway edge} dan \textit{cloud/PIP}. Komponen ini berfungsi sebagai perantara komunikasi antara \textit{gateway} dan \textit{cloud}. Dalam eksperimen, \textit{proxy/server} merepresentasikan jalur komunikasi menengah yang dilalui ketika \textit{gateway} harus mengambil atribut aktual dari \textit{cloud/PIP}. Dengan adanya \textit{proxy}, arsitektur dapat memodelkan adanya biaya komunikasi tambahan ketika evaluasi ABAC di \textit{edge} membutuhkan atribut yang tidak tersedia secara lokal.

Prinsip penting dalam arsitektur ini adalah bahwa seluruh atribut yang dibutuhkan oleh kebijakan ABAC tetap dievaluasi. \textit{Optimizer} tidak digunakan untuk menghapus atribut dari proses evaluasi kebijakan. Sebaliknya, \textit{optimizer} hanya menentukan atribut mana yang boleh di\textit{cache}, nilai TTL per atribut, dan ukuran \textit{cache} lokal pada \textit{gateway}. Dengan demikian, optimasi yang dilakukan tidak melemahkan semantik kebijakan ABAC, melainkan hanya mengatur strategi akuisisi atribut agar evaluasi akses menjadi lebih efisien tanpa mengabaikan kualitas keputusan otorisasi.

Alur evaluasi akses dimulai ketika \textit{endpoint} IoT mengirimkan permintaan akses ke PEP pada \textit{gateway edge}. PEP kemudian meneruskan permintaan tersebut ke PDP. PDP memeriksa atribut yang dibutuhkan melalui \textit{cache} atribut lokal. Jika atribut tersedia di \textit{cache} dan TTL-nya masih valid, maka nilai atribut dari \textit{cache} digunakan untuk evaluasi kebijakan. Kondisi ini disebut sebagai \textit{cache hit}. Sebaliknya, jika atribut tidak tersedia di \textit{cache} atau TTL telah kedaluwarsa, maka terjadi \textit{cache miss} atau \textit{TTL expired}. Pada kondisi tersebut, PDP mengirimkan \textit{attribute query} ke \textit{cloud/PIP} melalui \textit{proxy/server} untuk memperoleh nilai atribut aktual.

Setelah \textit{cloud/PIP} menerima \textit{query}, \textit{cloud/PIP} mengirimkan respons atribut kembali melalui \textit{proxy/server} menuju \textit{gateway}. Respons tersebut dapat berupa nilai atribut aktual, pembaruan atribut, atau informasi revokasi. Jika atribut tersebut termasuk dalam kategori \textit{cacheable attribute, gateway} dapat memperbarui \textit{cache} atribut lokal. Selanjutnya, PDP menggunakan atribut yang tersedia untuk mengevaluasi kebijakan ABAC dan menghasilkan \textit{candidate decision}. Keputusan tersebut diteruskan ke PEP untuk ditegakkan, lalu PEP mengirimkan respons keputusan berupa \textit{allow} atau \textit{deny} kepada \textit{endpoint}.

Arsitektur ini memungkinkan penelitian untuk mengevaluasi \textit{trade-off} antara performa dan keamanan. Dari sisi performa, penggunaan \textit{cache} lokal di \textit{gateway} dapat menurunkan altensi dan jumlah \textit{PIP calls}. Namun, dari sisi keamanan, penggunaan \textit{cache} juga dapat menimbulkan risiko \textit{stale attribute}, yaitu kondisi ketika nilai atribut di \textit{cache} tidak lagi sesuai dengan kondisi aktual. Oleh karena itu, penelitian ini tidak hanya mengukur metrik performa, tetapi juga metrik kualitas keputusan.

Secara konseptual, arsitektur sistem yang diusulkan dan alur komunikasi antar komponennya dapat digambarkan melalui \textit{sequence diagram} pada Gambar \ref{fig:seq-diagram-arsitektur}.

\vspace{0.5em}

\begin{figure}[H]
    \centering
    \includegraphics[width=1\textwidth]{seq-diagram-arsitektur.png}
    \caption{Sequence diagram arsitektur sistem yang diusulkan}
    \label{fig:seq-diagram-arsitektur}
\end{figure}

\vspace{-0.5em}

Gambar \ref{fig:seq-diagram-arsitektur} menunjukkan bahwa komunikasi antara \textit{gateway edge} dan \textit{cloud/PIP} bersifat dua arah. \textit{Gateway} tidak hanya mengirimkan \textit{attribute query}, tetapi juga menerima respons atribut, pembaruan atribut, atau pembaruan revokasi melalui \textit{proxy/server}. Setelah atribut diterima, PDP dapat memperbarui \textit{cache} untuk atribut yang dikonfigurasi sebagai \textit{cacheable attributes}. Selanjutnya PDP mengevaluasi kebijakan ABAC dan mengirimkan \textit{candidate decision} kepada PEP. PEP kemudian menegakkan keputusan tersebut dan mengembalikan respons keputusan kepada \textit{endpoint} IoT. Dengan alur ini, arsitektur yang diusulkan dapat mengevaluasi manfaat \textit{cache} atribut lokal sekaligus mengukur risiko keamanan akibat atribut yang tidak mutakhir.